\section{Proofs for original-law transport qualification}\label{supp:transport90}
\subsection{One-step comparison}
For probability measures $\mu,\nu$ and a $K$-Lipschitz function $f$, every coupling $\Gamma$ of the measures satisfies
\[
 \left|\int f\,d\mu-\int f\,d\nu\right|
 =\left|\int(f(y)-f(z))\,\Gamma(dy,dz)\right|
 \le K\int d(y,z)\,\Gamma(dy,dz).
\]
Taking an infimum over couplings proves the required $K W_1$ bound. The argument does not require densities or attainment of the infimum. Boundedness or the assumed finite first moments makes the integrals well defined.

Fix a qualified $\theta$ and suppress its subscript. Define the original-law action operator
\[
 Q_t^P f_{t+1}(x,a)=c_t(x,a)+\beta_tP_t^a f_{t+1}(x).
\]
The preceding coupling inequality and \eqref{eq:transport-input90} give, for every original feasible action,
\begin{equation}\label{eq:one-step90}
 |Q_t^P f_{t+1}(x,a)-\widehat Q_t(x,a)|\le b_t.
\end{equation}
If two functions on a nonempty action set differ uniformly by at most $b$, their infima differ by at most $b$: each function is everywhere at least the infimum of the other minus $b$, and the reverse argument gives the upper bound. No differentiability, uniqueness or continuity of the minimizing action is needed. It follows that
\begin{align}\label{eq:true-residual90}
 \ell_t-b_t&\le T_t f_{t+1}-f_t\le u_t+b_t,\notag\\
 \ell_t-b_t&\le T_t^\pi f_{t+1}-f_t
                         \le u_t+\gamma_t+b_t.
\end{align}
The lower policy inequality uses feasibility: the approximate action value at the implemented action is at least its infimum. The upper inequality uses the certified approximate greediness. All inequalities hold on the whole original state domain.

\subsection{Backward induction and the original optimum}
For a sequence $r_t$, define $S_t(r,d)=\sum_{s=t}^{T-1}(\prod_{j=t}^{s-1}\beta_j)r_s+(\prod_{j=t}^{T-1}\beta_j)d$. Empty products equal one. The claim
\begin{align*}
 S_t(\ell-b,d^-)&\le V_t^*-f_t\le S_t(u+b,d^+),\\
 S_t(\ell-b,d^-)&\le V_t^\pi-f_t
                            \le S_t(u+\gamma+b,d^+)
\end{align*}
follows by backward induction. At $T$, it is the terminal enclosure. Suppose the first inequality holds at $t+1$. Order preservation and the identity $T_t(v+k)=T_t v+\beta_t k$ give its lower bound at $t$ from the lower bound in \eqref{eq:true-residual90}. The upper bound follows in the same way. The fixed-policy operator has the same order and constant-shift properties, so the second line follows separately. This induction is under the original kernels; the numerical kernels do not determine the economic value being bounded.

Subtract the lower bound for $V_0^*-f_0$ from the upper bound for $V_0^\pi-f_0$. The result is exactly $\mathcal R_\theta$. Feasibility gives $V_0^\pi\ge V_0^*$ and proves \eqref{eq:transport-loss90}. The proof holds at each initial state and hence after averaging under any initial distribution supported on the declared domain. It also holds for every parameter for which the uniform hypotheses were certified. If the hypotheses are established on a single simultaneous data event, the conclusion holds on that event without an additional union over parameter vectors. Pointwise certificates obtained only at observed states do not create that simultaneous event.

The parameter may determine a different future continuation and a different future action at every date. We fixed it only to prove an inequality for its complete policy; the proof never substitutes another parameter's future. Discontinuous acquisition or deterministic tie-breaking therefore causes no problem as long as the resulting policy is measurable, feasible and satisfies the whole-domain residual inequalities. A transition bound applying only under the reference policy would not establish \eqref{eq:one-step90} and is insufficient.

\subsection{State-independent shifts and singular laws}
Replace $f_t$ by $f_t+k_t$ and $\widehat Q_t$ by $\widehat Q_t+\beta_t k_{t+1}$. Since both kernels have total mass one, the arithmetic discrepancy, transport modulus and greedy gap are unchanged. The approximate Bellman residual is shifted by $\beta_t k_{t+1}-k_t$. Its certified lower and upper endpoints shift by that same amount, so their difference is unchanged. The terminal enclosure endpoints both decrease by $k_T$. Every term of $\mathcal R_\theta$ is consequently invariant. Using absolute levels of the continuation in place of its oscillation would forfeit this property.

If $\widehat P_t$ is the law of $q_t(Y)$ when $Y$ has law $P_t$, the joint law of $(Y,q_t(Y))$ is a coupling. Its expected distance is at most $h_t$. Thus $W_1(P_t,\widehat P_t)\le h_t$. Substitution into $2\sum_tB_tb_t$ gives \eqref{eq:transport-quantization90}. In particular, $\delta_h$ and $\delta_0$ have transport distance $h$ but total-variation distance $\sup_A|\delta_h(A)-\delta_0(A)|=1$ when $h>0$. Mutual singularity therefore does not obstruct this continuation-based calculation. The conclusion concerns the transport contribution only; a nonvanishing residual span or greedy error still has to be paid.

\subsection{Stationary discounted bound}
Let $T$ and $T^\pi$ be the bounded discounted Bellman and policy operators. They preserve order, translate constants by $\beta$ and are sup-norm contractions with modulus $\beta$. Iterating \eqref{eq:true-residual90} from the bounded function $f$, geometric summation yields
\[
 \frac{\ell-b}{1-\beta}\le V^*-f\le\frac{u+b}{1-\beta},
 \qquad
 V^\pi-f\le\frac{u+\gamma+b}{1-\beta}.
\]
The contraction limits are the original optimal and policy values. Alternatively, apply the finite-horizon induction and let the bounded discounted terminal difference tend to zero. Subtraction proves Corollary~\ref{cor:stationary90}. This argument requires a certified modulus of the stored $f$; it does not infer a uniform modulus for every optimal value from Lipschitz primitives. That distinction matters when transition dynamics expand distances.

\subsection{Resource allocation and the other economic objectives}
With $\eta_t\ge u_t-\ell_t+\gamma_t$, the sufficient budget \eqref{eq:joint-budget90} is obtained by rearranging the nonnegative terms in $\mathcal R_\theta$. After subtracting fixed terminal and arithmetic contributions, it is a linear budget in the residual and transport allowances. On positive-weight coordinates, any resource upper model satisfying Theorem~\ref{thm:coupled89} has that theorem's convexity, attainment, dyadic-rounding and dual-certificate properties. This is a direct application with additional accuracy coordinates, not a newly asserted lower bound on all numerical methods. Caps and certificates for the moduli are part of the input. Optimizing a modulus jointly requires retaining its parameter dependence in the resource and error constraints.

For a cash-invariant monotone one-step risk operator with discount $\beta_t$, a certified numerical discrepancy $b_t$ gives the same two residual inequalities and the same constant-shift induction. Neither the coupling proof for linear expectation nor cash invariance alone supplies that numerical discrepancy for a nonlinear risk functional. For a general Lipschitz aggregator, the separate absolute-error induction in Proposition~\ref{prop:nonlinear86} supplies the valid replacement. Applying an original-optimum bound to each fixed-opponent best-response problem gives the stated unilateral incentive account; it does not prove existence of a game equilibrium. These arguments preserve the distinct hypotheses of the original applications.

\subsection{Diagnostic with a falsely free transition approximation}
In the one-date example, both numerical action values are constants: zero for action $0$ and $h/2$ for action $1$. Hence $\widehat T f_1-f_0=0$, the numerical greedy gap is zero, and the terminal enclosure is exact. The true values of the two actions are $h$ and $h/2$, respectively. The numerical greedy policy therefore has regret $h/2$. The stored terminal function has modulus one and the largest transport discrepancy is $h$, so Theorem~\ref{thm:transport90} gives $2h$. The omitted-charge expression is zero and is invalid. All assertions are exact for rational $h$; no sampling or floating-point limit is used.
